% ============================================================
% sialkot_football_marxian_example.tex
% Standalone LaTeX file (pdflatex) — no external dependencies.
% ============================================================

\documentclass[11pt]{article}

\usepackage[a4paper,margin=1in]{geometry}
\usepackage{amsmath,amssymb,mathtools}
\usepackage{booktabs,longtable}
\usepackage{siunitx}
\usepackage{enumitem}
\usepackage{microtype}
\usepackage[hidelinks]{hyperref}

\sisetup{
  detect-all,
  group-separator={,},
  group-minimum-digits=4
}

% ---------- Convenience macros ----------
\newcommand{\PKR}{\text{PKR}}
\newcommand{\USD}{\text{USD}}
\newcommand{\uBall}{\text{ball}}
\newcommand{\uHour}{\text{hour}}
\newcommand{\uDay}{\text{day}}
\newcommand{\uMonth}{\text{month}}
\newcommand{\uYear}{\text{year}}

\newcommand{\cconst}{c} % constant capital
\newcommand{\vvar}{v}   % variable capital
\newcommand{\splus}{s}  % surplus value (here proxied by accounting profit)
\newcommand{\kcost}{k}  % cost price: k=c+v
\newcommand{\eRate}{e}  % rate of exploitation: e=s/v
\newcommand{\pRate}{r}  % profit rate: r=s/(c+v)
\newcommand{\OCC}{\mathrm{OCC}} % organic composition proxy: c/v

\title{A Worked, Real-World Marxian Example\\Football Manufacturing in Sialkot (Pakistan)}
\author{Standalone teaching note (with fully worked arithmetic)}
\date{\today}

% Figures shared with the pamphlet and slides, generated from data/sialkot.json
% Written by tools/sialkot.py from data/sialkot.json. Do not edit by hand.
\newcommand{\PieceRatePKR}{160}
\newcommand{\BallsPerDay}{3}
\newcommand{\HoursPerBall}{3}
\newcommand{\WorkDaysMonth}{20}
\newcommand{\LivingWageQuotePKR}{20000}
\newcommand{\LivingWageUrbanPKR}{51733}
\newcommand{\LivingWageRuralPKR}{45588}
\newcommand{\cPKR}{154.52}
\newcommand{\vPKR}{56.29}
\newcommand{\kPKR}{210.83}
\newcommand{\MarkupPct}{8.1}
\newcommand{\pPKR}{227.91}
\newcommand{\sPKR}{17.08}
\newcommand{\NewValuePKR}{73.37}
\newcommand{\RetailUSD}{165}
\newcommand{\ExportAvgUSD}{4.44}


\begin{document}
\maketitle

\vspace{-0.5em}
\begin{center}
\small
This document uses real, published numbers on Sialkot football manufacturing to build two concrete worked examples:
(1) a hand-stitcher paid per ball; and (2) the per-ball cost structure of a Sialkot firm.
Every calculation is shown step-by-step and explained in plain language.
\end{center}

\section{Industry snapshot (why Sialkot matters)}
Sialkot is a globally important production cluster for footballs. Bloomberg reports that the city produces about
\emph{70\% of the world’s soccer balls}, with roughly \emph{1,000 factories} and about \emph{60,000 people}
working in the soccer ball manufacturing business.\footnote{Bloomberg Businessweek, ``This Is Where Most of the World’s Soccer Balls Come From'' (Nov 22, 2022). \url{https://www.bloomberg.com/features/2022-world-cup-soccer-ball-adidas-al-rihla-sialkot/}}
Bloomberg also reports that \emph{more than 80\%} of Sialkot’s soccer balls use \emph{hand stitching}, though thermo-bonding
is expanding.\footnote{Same Bloomberg source as above.}

On exports: Pakistan’s trade reporting (TDAP) lists football exports of about \USD{}~254.0 million in FY2023--24 (and \USD{}~237.0 million in FY2022--23).\footnote{TDAP, \emph{Pakistan Trade Perspective} FY2023--24 and TDAP annual trade performance FY2022--23 (tables listing ``Footballs''). \url{https://tdap.gov.pk/wp-content/uploads/2025/01/TDAP-Annual-Report-PAKISTAN-TRADE-PERSPECTIVE-FY-23-24-new.pdf} and \url{https://tdap.gov.pk/wp-content/uploads/2023/10/Final-draft-of-Pakistans-Annual-Trade-Performance-July-June-FY23-6-SEPTEMBER-2023.pdf}}

A Reuters report (reprinted by \emph{Dawn}) adds two useful scale facts:
Forward Sports said it made \emph{5.5 million} Al Rihla balls, and the local chamber claimed
\emph{43 million balls valued at \USD{}~191 million} were made in Sialkot in FY2021--22.\footnote{\emph{Dawn} (Reuters), ``‘Made-in-Sialkot’ ball puts Pakistan in FIFA World Cup'' (Dec 9, 2022). \url{https://www.dawn.com/news/1725468}}

\section{What we will compute (and why the steps are so slow)}
We will compute three things, carefully:
\begin{enumerate}[leftmargin=*]
\item \textbf{Piece-rate wages in time form.} A ``per ball'' payment is just a wage in disguise. If we know how many hours it takes to stitch one ball, we can compute an hourly wage.
\item \textbf{Per-ball cost structure.} Using firm microdata, we will split the cost of one ball into \emph{materials}, \emph{overhead}, and \emph{wages}.
\item \textbf{Marxian ratios in money form.} With those components, we compute:
  \[
    \kcost=\cconst+\vvar,\qquad
    \splus,\qquad
    \eRate=\frac{\splus}{\vvar},\qquad
    \pRate=\frac{\splus}{\cconst+\vvar},\qquad
    \OCC=\frac{\cconst}{\vvar}.
  \]
\end{enumerate}

\section{Mini arithmetic primer (only what we need)}
This section is intentionally basic.

\subsection*{(A) Multiplication}
If you earn \(\PKR{}~160\) per ball and you stitch \(3\) balls:
\[
  160 \times 3 = 480.
\]
That means \(\PKR{}~480\) total.

\subsection*{(B) Division}
If you earned \(\PKR{}~480\) over \(9\) hours:
\[
  \frac{480}{9} \approx 53.33.
\]
This is ``\(\PKR{}~53.33\) per hour''.

\subsection*{(C) Turning a ratio into a percent}
If a number is \(0.48\) and you want percent:
\[
  0.48 \times 100 = 48\%.
\]

\section{Notation (definitions you can check)}
We use standard Marxian accounting notation:
\begin{itemize}[leftmargin=*]
\item \(\cconst\): \textbf{constant capital} (money spent on materials, energy, depreciation---inputs that transfer value to output).
\item \(\vvar\): \textbf{variable capital} (wages: money spent on labour-power).
\item \(\kcost\): \textbf{cost price}, defined as \(\kcost=\cconst+\vvar\).
\item \(\splus\): \textbf{surplus value}. In this document we will \emph{approximate} \(\splus\) by observed accounting profit at the producing firm level, for a transparent first pass.
\item \(\eRate\): \textbf{rate of exploitation}, \(\eRate=\splus/\vvar\).
\item \(\pRate\): \textbf{profit rate}, \(\pRate=\splus/(\cconst+\vvar)\).
\item \(\OCC\): \textbf{organic composition proxy}, \(\OCC=\cconst/\vvar\).
\end{itemize}

\bigskip
\noindent
\textbf{Important limitation (stated up front):}
a low observed producer profit margin does \emph{not} imply low exploitation.
It can also reflect intense competition, buyer power, and value capture downstream (brands, logistics, retail).
We will show this tension explicitly.

% ============================================================
\section{Worked Example 1: a hand-stitcher paid per ball (Bloomberg data)}
% ============================================================

Bloomberg reports the following for stitchers at manufacturer Anwar Khawaja Industries:
\begin{itemize}[leftmargin=*]
\item payment per ball: \(w_{\uBall}=\PKR{}~160\) (about \USD{}~0.75),
\item time per ball: \(t_{\uBall}=3\) hours,
\item output per day (typical): \(q_{\uDay}=3\) balls,
\item monthly earnings at that pace: about \(\PKR{}~9{,}600\),
\item an estimated living wage for Sialkot: about \(\PKR{}~20{,}000\) per month.\footnote{Bloomberg Businessweek (Nov 22, 2022), as cited earlier.}
\end{itemize}

\subsection{Step 1: daily earnings from a piece rate}
Daily earnings are:
\[
  w_{\uDay} = w_{\uBall}\times q_{\uDay}.
\]
Now substitute the numbers:
\[
  w_{\uDay} = 160 \times 3.
\]
Compute \(160\times 3\):
\[
  160\times 3 = 480.
\]
So:
\[
  w_{\uDay}=\PKR{}~480.
\]

\subsection{Step 2: implied working hours per day}
If each ball takes \(3\) hours, and the stitcher completes \(3\) balls:
\[
  T_{\uDay} = t_{\uBall}\times q_{\uDay} = 3\times 3 = 9~\text{hours}.
\]
So the stitching time alone is about \(9\) hours per day.

\subsection{Step 3: hourly wage implied by the piece rate}
Hourly wage is:
\[
  w_{\uHour} = \frac{w_{\uDay}}{T_{\uDay}}.
\]
Substitute:
\[
  w_{\uHour}=\frac{480}{9}.
\]
Do the division:
\[
  \frac{480}{9}\approx 53.33.
\]
So:
\[
  w_{\uHour}\approx \PKR{}~53.33/\uHour.
\]

\subsection{Step 4: monthly earnings (showing the hidden assumption)}
Bloomberg reports monthly earnings of about \(\PKR{}~9{,}600\) at this pace.
We can recover the implied number of working days.

We already have \(w_{\uDay}=\PKR{}~480\).
Let the number of working days per month be \(D_{\uMonth}\).
Then:
\[
  w_{\uMonth}=w_{\uDay}\times D_{\uMonth}.
\]
We set \(w_{\uMonth}=\PKR{}~9{,}600\) and solve for \(D_{\uMonth}\):
\[
  9{,}600 = 480\times D_{\uMonth}.
\]
Divide both sides by \(480\):
\[
  D_{\uMonth}=\frac{9{,}600}{480}.
\]
Now compute:
\[
  \frac{9{,}600}{480} = 20.
\]
So the implied working schedule is about:
\[
  D_{\uMonth}\approx 20~\uDay/\uMonth.
\]

\subsection{Step 5: compare to a living-wage benchmark}
Living wage benchmark:
\[
  L_{\uMonth}=\PKR{}~20{,}000/\uMonth.
\]
Actual earnings:
\[
  w_{\uMonth}=\PKR{}~9{,}600/\uMonth.
\]
Compute the fraction:
\[
  \frac{w_{\uMonth}}{L_{\uMonth}}=\frac{9{,}600}{20{,}000}.
\]
Divide numerator and denominator by \(100\) to simplify:
\[
  \frac{9{,}600}{20{,}000}=\frac{96}{200}.
\]
Now divide:
\[
  \frac{96}{200}=0.48.
\]
Convert to percent:
\[
  0.48\times 100 = 48\%.
\]
So the reported monthly earnings are about:
\[
  48\% \text{ of the cited living-wage benchmark.}
\]

\subsection{Step 6: share of the retail price (Al Rihla launch price)}
A contemporaneous report states the Al Rihla official match ball retailed at \USD{}~165 on Adidas’ US website at launch.\footnote{GOAL (Mar 30, 2022), retail price table. \url{https://www.goal.com/en/news/adidas-world-cup-match-ball-al-rihla-qatar-2022/blta7dc0012ef2a879f}}

Bloomberg reports stitchers get about \USD{}~0.75 per ball.
Let the stitcher’s pay per ball be \(w^{\$}_{\uBall}=0.75\) and the retail price be \(P^{\$}_{\text{retail}}=165\).

The share is:
\[
  \text{share}=\frac{w^{\$}_{\uBall}}{P^{\$}_{\text{retail}}}=\frac{0.75}{165}.
\]
Compute the division:
\[
  \frac{0.75}{165}\approx 0.004545\ldots
\]
Convert to percent:
\[
  0.004545\ldots \times 100 \approx 0.4545\%.
\]
So, at launch pricing, the stitcher’s pay is roughly:
\[
  0.45\% \text{ of the retail price of an Al Rihla match ball.}
\]

\subsection{Step 7: wage-hours needed to buy the commodity}
How long must the stitcher work to buy back, at the retail benchmark, the kind of commodity this labour helps to produce?

First put the retail price into rupees. Step 6 used Bloomberg's dollar figure for pay, from 2022. Here we use one stated reference rate instead, \(\PKR{}~\USDPKR\) per \USD{}~1 (NBP Treasury rate sheet, \RateDate), so the two steps should not be read as a single conversion:
\[
  P_{\text{retail}} = 165\times \USDPKR \approx \PKR{}~\SkRetail.
\]
Now divide by the hourly wage from Step 3. Call the result \(H\), the wage-hours needed:
\[
  H=\frac{P_{\text{retail}}}{w_{\uHour}}\approx\frac{\SkRetail}{\SkHourly}\approx \SkHours\ \text{hours}.
\]
At \(3\) hours per ball, that is the time taken to stitch:
\[
  \frac{\SkHours}{3}\approx \SkBalls\ \text{balls}.
\]
At \(9\) hours per day, it is:
\[
  \frac{\SkHours}{9}\approx \SkDays\ \text{workdays}.
\]
So, at these benchmarks, a worker would stitch about \SkBalls{} balls, over about \SkDays{} nine-hour days, to earn the retail price of one.

\textbf{How to read this.} The benchmark is the official Al Rihla match ball, a premium product. So this measures the distance between the wage and that market price, and not the price of the particular ball this worker stitched. It also counts all of the wage towards the purchase, leaving nothing for food, rent or anything else.

\subsection{A quick ``order'' calculation (1,000 balls)}
Suppose a buyer places an order for \(Q=1{,}000\) hand-stitched balls at \(160\) rupees stitching pay each.

Total stitching wage bill for that order:
\[
  W_{\text{stitch}} = 160\times 1{,}000.
\]
Compute:
\[
  160\times 1{,}000 = 160{,}000.
\]
So:
\[
  W_{\text{stitch}}=\PKR{}~160{,}000.
\]

If each ball takes \(3\) hours, total stitching hours:
\[
  H_{\text{stitch}} = 3\times 1{,}000 = 3{,}000~\text{hours}.
\]
This is useful later: it lets you translate piece-wages into labour-time at the batch scale.

% ============================================================
\section{Worked Example 2: per-ball cost structure (Atkin et al.\ microdata)}
% ============================================================

Now we move from one worker’s wage to one firm’s \emph{per-ball} cost structure using microdata
from Atkin et al.\ (NBER), based on a baseline survey of Pakistani soccer ball producers.\footnote{Atkin, Khandelwal, Osman, and Pavcnik (2015), NBER Working Paper 21112. \url{https://www.nber.org/papers/w21112}}

\subsection{The raw data (one size-5 promotional ball)}
Atkin et al.\ report an input cost table (``Table A.1: input cost in rupees'') for a size-5, two-layer promotional ball
and note an exchange rate of about \(\PKR{}~100=\USD{}~1\) for their tabulation.\footnote{Atkin et al.\ (2015), Appendix Table A.1 and note on exchange rate. \url{https://www.nber.org/papers/w21112}}

Here are the per-ball money costs (rupees):
\begin{center}
\begin{longtable}{@{}lS[table-format=3.2]@{}}
\toprule
\textbf{Cost item (per ball)} & {\textbf{Rupees}} \\
\midrule
Rexine (outer synthetic material)      & 39.68 \\
Backing cloth                          & 23.27 \\
Latex                                  & 38.71 \\
Bladder                                & 42.02 \\
Cutting labour                          & 1.49 \\
Stitching labour                        & 39.24 \\
Other labour                            & 15.56 \\
Overhead                                & 10.84 \\
\midrule
\textbf{Total cost (all items)}         & 210.83 \\
\bottomrule
\end{longtable}
\end{center}

\subsection{Step 1: classify items into \(\cconst\) and \(\vvar\)}
We treat \textbf{materials} and \textbf{overhead} as constant capital \(\cconst\),
and \textbf{labour} items as variable capital \(\vvar\).

So:

\paragraph{Constant capital \(\cconst\) (materials + overhead).}
\[
  \cconst = 39.68 + 23.27 + 38.71 + 42.02 + 10.84.
\]
Add step by step:
\[
  39.68 + 23.27 = 62.95
\]
\[
  62.95 + 38.71 = 101.66
\]
\[
  101.66 + 42.02 = 143.68
\]
\[
  143.68 + 10.84 = 154.52
\]
So:
\[
  \cconst = \PKR{}~154.52.
\]

\paragraph{Variable capital \(\vvar\) (wages).}
\[
  \vvar = 1.49 + 39.24 + 15.56.
\]
Add step by step:
\[
  1.49 + 39.24 = 40.73
\]
\[
  40.73 + 15.56 = 56.29
\]
So:
\[
  \vvar = \PKR{}~56.29.
\]

\paragraph{Check: \(\kcost=\cconst+\vvar\).}
\[
  \kcost=\cconst+\vvar=154.52+56.29=210.81.
\]
This is (up to rounding of cents) the reported total \(210.83\). Small differences like \(0.02\) happen because
the table entries are rounded to two decimals.

\subsection{Step 2: a profit margin (what Atkin et al.\ report)}
Atkin et al.\ also report profit margins; for size-5 promotional balls, a mean ``profit (\%)'' around \(8.1\%\) is shown in their table.\footnote{Atkin et al.\ (2015), Table of average prices and profit percentages for ball types (e.g.\ size-5 promotional). \url{https://www.nber.org/papers/w21112}}
They also state the conversion they use: \emph{unit cost} is recovered by dividing price by \(1+\) (profit margin),
which means the margin is a \textbf{markup on cost}.\footnote{Atkin et al.\ (2015), note describing ``divide its reported price by one plus its reported profit margin'' to obtain unit cost. \url{https://www.nber.org/papers/w21112}}

So we will treat:
\[
  m = 0.081
\]
as the (mean) profit markup on cost for this illustrative ball.

\subsection{Step 3: compute an implied producer selling price}
If markup on cost is \(m\), then:
\[
  p = \kcost(1+m).
\]
Use \(\kcost=210.83\) and \(m=0.081\):
\[
  p = 210.83(1+0.081)=210.83\times 1.081.
\]

Now compute \(210.83\times 1.081\) carefully by splitting \(1.081\):
\[
  1.081 = 1 + 0.081.
\]
So:
\[
  210.83\times 1.081 = 210.83\times 1 + 210.83\times 0.081.
\]
The first term is just \(210.83\).

Now compute the second term:
\[
  210.83\times 0.081 = 210.83\times \frac{81}{1000}.
\]
First compute \(210.83\times 81\):
\[
  210.83\times 81 = 210.83\times 80 + 210.83\times 1.
\]
\[
  210.83\times 80 = (210.83\times 8)\times 10.
\]
Compute \(210.83\times 8\):
\[
  210.83\times 8 = 1686.64.
\]
So:
\[
  210.83\times 80 = 1686.64\times 10 = 16866.40.
\]
Add the \(\times 1\) part:
\[
  210.83\times 81 = 16866.40 + 210.83 = 17077.23.
\]
Now divide by \(1000\):
\[
  210.83\times 0.081 = \frac{17077.23}{1000} = 17.07723.
\]

So the price is:
\[
  p = 210.83 + 17.07723 = 227.90723.
\]
Rounded to two decimals:
\[
  p \approx \PKR{}~227.91 \quad \text{(producer selling price per ball).}
\]

\subsection{Step 4: compute \(\splus\) as the per-ball accounting profit}
Accounting profit per ball is:
\[
  \pi = p - \kcost.
\]
Substitute:
\[
  \pi = 227.90723 - 210.83 = 17.07723.
\]
In this simplified first pass, we set:
\[
  \splus \approx 17.08.
\]

\subsection{Step 5: compute the Marxian ratios}
Now we compute:
\[
  \eRate=\frac{\splus}{\vvar},\quad
  \pRate=\frac{\splus}{\cconst+\vvar},\quad
  \OCC=\frac{\cconst}{\vvar}.
\]

\paragraph{(A) Rate of exploitation \(\eRate=\splus/\vvar\).}
\[
  \eRate = \frac{17.07723}{56.29}.
\]
A quick long-division approximation:
\[
  56.29\times 0.30 = 16.887 \quad (\text{slightly below }17.07723),
\]
difference:
\[
  17.07723 - 16.887 = 0.19023.
\]
Now add \(0.003\) more:
\[
  56.29\times 0.003 = 0.16887,
\]
new difference:
\[
  0.19023 - 0.16887 = 0.02136.
\]
Add \(0.0003\) more:
\[
  56.29\times 0.0003 = 0.016887,
\]
difference:
\[
  0.02136 - 0.016887 = 0.004473.
\]
So we are at \(0.3033\) with a tiny remainder. Therefore:
\[
  \eRate \approx 0.303 \approx 30.3\%.
\]

\paragraph{(B) Profit rate \(\pRate=\splus/(\cconst+\vvar)\).}
Here \(\cconst+\vvar \approx 210.83\), so:
\[
  \pRate = \frac{17.07723}{210.83} \approx 0.081 = 8.1\%.
\]
(This matches the assumed markup construction.)

\paragraph{(C) Organic composition proxy \(\OCC=\cconst/\vvar\).}
\[
  \OCC = \frac{154.52}{56.29}.
\]
Compute:
\[
  56.29\times 2 = 112.58,\quad \text{remainder } 154.52-112.58=41.94.
\]
Now:
\[
  56.29\times 0.7 = 39.403,\quad \text{remainder } 41.94-39.403=2.537.
\]
Now:
\[
  56.29\times 0.04 = 2.2516,\quad \text{remainder } 2.537-2.2516=0.2854.
\]
Now:
\[
  56.29\times 0.005 = 0.28145,\quad \text{remainder } 0.2854-0.28145\approx 0.00395.
\]
So:
\[
  \OCC \approx 2.745 \approx 2.75.
\]

\subsection{Step 6: ``necessary'' vs ``surplus'' shares (labour-time intuition without MELT)}
Even if you do not know the MELT, one identity is helpful.

New value (in money terms) created in production is:
\[
  \text{new value} = \vvar + \splus.
\]
Here:
\[
  \vvar + \splus \approx 56.29 + 17.08 = 73.37.
\]
Wage share of new value:
\[
  \frac{\vvar}{\vvar+\splus} = \frac{56.29}{73.37}\approx 0.767 = 76.7\%.
\]
Surplus share:
\[
  \frac{\splus}{\vvar+\splus} = \frac{17.08}{73.37}\approx 0.233 = 23.3\%.
\]
In labour-time language, these shares correspond to:
\[
  \frac{T_n}{T}=\frac{1}{1+\eRate},\qquad
  \frac{T_s}{T}=\frac{\eRate}{1+\eRate}.
\]
With \(\eRate\approx 0.303\):
\[
  \frac{T_n}{T}\approx \frac{1}{1.303}\approx 0.767,\quad
  \frac{T_s}{T}\approx \frac{0.303}{1.303}\approx 0.233.
\]
So, \emph{if} the relevant working day were, say, \(T=9\) hours, then:
\[
  T_n \approx 0.767\times 9 \approx 6.90~\uHour,\qquad
  T_s \approx 0.233\times 9 \approx 2.10~\text{hours}.
\]
That is: roughly \(6.9\) hours to reproduce wages, and \(2.1\) hours as surplus labour time---under the simplifying
assumption that the money shares reflect value shares and that the ratios apply uniformly across the direct labour in this production process.

\section{Scaling up: from a ball to the Sialkot cluster (transparent approximations)}
Here we take published scale facts and do simple ``back-of-the-envelope'' totals. The point is not precision; the point is
to show exactly what assumptions you must make, and how each assumption changes the result.

\subsection{Step 1: average export value per ball (Reuters/Dawn numbers)}
Reuters (via \emph{Dawn}) reports a chamber claim: \(43\) million balls valued at \USD{}~191 million (FY2021--22).\footnote{\emph{Dawn} (Reuters) (Dec 9, 2022), as cited earlier.}

Average value per ball is:
\[
  \overline{P} = \frac{191{,}000{,}000}{43{,}000{,}000}.
\]
Cancel six zeros by dividing numerator and denominator by \(1{,}000{,}000\):
\[
  \overline{P} = \frac{191}{43}.
\]
Now divide \(191\) by \(43\):
\[
  43\times 4 = 172,\quad \text{remainder } 191-172 = 19.
\]
So:
\[
  \frac{191}{43} = 4 + \frac{19}{43}.
\]
And:
\[
  \frac{19}{43}\approx 0.44186.
\]
So:
\[
  \overline{P}\approx 4.44~\USD/\uBall.
\]
Interpretation: the average value per ball at that scale is about \USD{}~4.44.

\subsection{Step 2: total stitching wage bill under two piece-rate scenarios}
We will compute two \emph{stitching-only} wage-bills:

\paragraph{Scenario A (high-end hand stitching; Bloomberg).}
Bloomberg reports about \USD{}~0.75 stitching pay per ball.\footnote{Bloomberg Businessweek (Nov 22, 2022), as cited earlier.}
Then stitching wage bill for \(Q=43\) million balls is:
\[
  W^{A}_{\text{stitch}} = 0.75\times 43{,}000{,}000.
\]
Compute:
\[
  0.75\times 43{,}000{,}000 = \frac{3}{4}\times 43{,}000{,}000 = 32{,}250{,}000.
\]
So:
\[
  W^{A}_{\text{stitch}}\approx 32.25~\text{million USD}.
\]

\paragraph{Scenario B (promotional ball stitching cost; Atkin et al.).}
Atkin et al.\ report stitching labour cost \(\PKR{}~39.24\) per ball, and use \(\PKR{}~100=\USD{}~1\) for their tabulation.\footnote{Atkin et al.\ (2015), Appendix Table A.1 and exchange-rate note.}
So stitching labour cost in dollars (for their dataset) is approximately:
\[
  w^{B}_{\uBall} \approx \frac{39.24}{100} = 0.3924~\USD/\uBall.
\]
Now multiply by \(43\) million balls:
\[
  W^{B}_{\text{stitch}} = 0.3924\times 43{,}000{,}000.
\]
Compute \(0.3924\times 43\) first:
\[
  0.3924\times 40 = 15.696,\qquad 0.3924\times 3 = 1.1772.
\]
Add:
\[
  15.696 + 1.1772 = 16.8732.
\]
So:
\[
  W^{B}_{\text{stitch}} \approx 16.8732~\text{million USD}.
\]

\subsection{Step 3: compare stitching wages to the export value}
Export value (Reuters/Dawn chamber claim): \USD{}~191 million.

\paragraph{Scenario A share.}
\[
  \text{share}^{A}=\frac{32.25}{191}\approx 0.1689 = 16.9\%.
\]

\paragraph{Scenario B share.}
\[
  \text{share}^{B}=\frac{16.87}{191}\approx 0.0883 = 8.8\%.
\]

\noindent
So stitching wages alone are somewhere between about \(9\%\) and \(17\%\) of export value \emph{under these two
very different ball types / wage regimes}.
This is not a contradiction: it reflects product heterogeneity (cheap vs premium balls), technology choice (machine vs hand stitching),
and buyer-driven price pressure.

\subsection{A note on why producer profit margins can be low while retail prices are high}
Atkin et al.\ report producer profit margins on the order of single digits (around \(8\%\) in the example table),
which is compatible with fierce competition among suppliers.\footnote{Atkin et al.\ (2015), profit table, as cited earlier.}
Meanwhile, retail pricing for the branded match ball can be on the order of \USD{}~165.\footnote{GOAL (Mar 30, 2022), as cited earlier.}
That wedge is where downstream costs and profits sit: branding, sponsorship, distribution, retail markups, and (often) monopoly rents.
This document does not ``solve'' that wedge; it makes it visible and computable.

\section{Historical note (labour regulation and the child-labour scandal)}
For context: the Sialkot cluster became a major site of international scrutiny over child labour in the 1990s,
followed by monitoring and restructuring initiatives.\footnote{For one overview: UNIDO (2010), Khalid Nadvi, ``The Soccer Ball Cluster of Sialkot, Pakistan'' (PDF). \url{https://www.unido.org/sites/default/files/files/2010-12/GC.13_10_Ebook_part2.pdf}}
This history matters because it shaped subcontracting patterns, home-based work, and compliance regimes that still structure
how piece-rate labour is organised.

\section{What you can change (and recompute immediately)}
If you want to adapt this to a specific pamphlet chapter, or to a different ball type, change these inputs and redo the same steps:
\begin{itemize}[leftmargin=*]
\item piece rate per ball (\(w_{\uBall}\)),
\item hours per ball (\(t_{\uBall}\)),
\item balls per day (\(q_{\uDay}\)),
\item cost table (materials, overhead, wage items),
\item markup / profit margin (\(m\)),
\item annual output volume (\(Q\)) and export value (for scaling).
\end{itemize}

\section*{References (quick list)}
\begin{thebibliography}{9}
\setlength{\itemsep}{0.35em}
\bibitem{Bloomberg2022}
Bloomberg Businessweek (2022).
\newblock ``This Is Where Most of the World’s Soccer Balls Come From'' (Nov 22, 2022).
\newblock \url{https://www.bloomberg.com/features/2022-world-cup-soccer-ball-adidas-al-rihla-sialkot/}

\bibitem{TDAPFY2324}
Trade Development Authority of Pakistan (TDAP) (2025).
\newblock \emph{Pakistan Trade Perspective FY 2023--24}.
\newblock \url{https://tdap.gov.pk/wp-content/uploads/2025/01/TDAP-Annual-Report-PAKISTAN-TRADE-PERSPECTIVE-FY-23-24-new.pdf}

\bibitem{TDAPFY23}
Trade Development Authority of Pakistan (TDAP) (2023).
\newblock \emph{Pakistan’s Annual Trade Performance July--June FY 2022--23}.
\newblock \url{https://tdap.gov.pk/wp-content/uploads/2023/10/Final-draft-of-Pakistans-Annual-Trade-Performance-July-June-FY23-6-SEPTEMBER-2023.pdf}

\bibitem{Atkin2015}
Atkin, D., Khandelwal, A.\ K., Osman, A., \& Pavcnik, N. (2015).
\newblock ``Exporting and Firm Performance: Evidence from a Randomized Experiment.''
NBER Working Paper 21112.
\newblock \url{https://www.nber.org/papers/w21112}

\bibitem{DawnReuters2022}
\emph{Dawn} (Reuters) (2022).
\newblock ``‘Made-in-Sialkot’ ball puts Pakistan in FIFA World Cup'' (Dec 9, 2022).
\newblock \url{https://www.dawn.com/news/1725468}

\bibitem{GoalPrice2022}
GOAL (2022).
\newblock ``Adidas unveils 'fastest in flight' World Cup match ball Al Rihla'' (Mar 30, 2022).
\newblock \url{https://www.goal.com/en/news/adidas-world-cup-match-ball-al-rihla-qatar-2022/blta7dc0012ef2a879f}

\bibitem{NadviUNIDO2010}
UNIDO (2010).
\newblock Khalid Nadvi, ``The Soccer Ball Cluster of Sialkot, Pakistan'' (PDF).
\newblock \url{https://www.unido.org/sites/default/files/files/2010-12/GC.13_10_Ebook_part2.pdf}
\end{thebibliography}

\end{document}
